\section{Released dictionaries are not near the floor}
\label{sec:dictionaries}

A sparse autoencoder's decoder is an overcomplete code in exactly the sense of
Section~\ref{sec:setting}: $F$ unit-norm directions in $d<F$ dimensions, used in practice by taking
inner products with feature directions. It is therefore an object the diagnostic applies to without
adaptation, and it is one that people deploy. This section measures \SaeDicts{} of them.

\paragraph{What is and is not computed.} Everything below comes from a decoder matrix. No model
weights are loaded, no activations are collected, no autoencoder is trained, and no linear programme
is solved. The cost is seconds per dictionary, which is why the section can span five model families
rather than one. The price is that these are statements about the geometry a dictionary presents to a
linear reader, not about what its features mean or how well it reconstructs; the reconstruction
metrics its authors report are a different and complementary measurement.

\paragraph{Design.} Four axes, chosen so that each isolates one variable
(Table~\ref{tab:sae}). \emph{Scale}: Qwen-Scope residual dictionaries at $d=2048$ and $d=4096$ with
load and recipe held at $16\times$ and TopK. \emph{Load}: EleutherAI's Pythia-160m dictionaries at
one layer with $k$ fixed at 64 and width varying, giving \SaeLoadOneLoad$\times$ to
\SaeLoadThreeLoad$\times$ at constant $d$. \emph{Causal}: EleutherAI's SmolLM2-135M pair, one arm
trained on the model and one on a randomly initialised copy of it, same recipe, same data, same ten
layers. \emph{Architecture}: the MLP down-projection of LFM2.5-350M, a hybrid convolution/attention
model rather than a standard transformer.

Every code is compared against an i.i.d.\ Gaussian code of \emph{its own} shape, evaluated at
\emph{its own} published operating sparsity. That matching is not cosmetic. Because $\sum_i h_i=d$,
the mean leverage of any code is exactly $d/F$, so absolute quantities like the fraction of features
a bound rules out are properties of the shape and not of the dictionary. Only the departure from a
shape-matched control carries information about the code.

\subsection{Near-floor geometry is generic}

The \SaeControls{} matched controls attain the floor: $\Rgeom$ from \SaeCtlRgeomMin{} to
\SaeCtlRgeomMax{} across every shape and every load from \SaeLoadMin$\times$ to
\SaeLoadMax$\times$, with leverage dispersion \SaeCtlCvMin{} to \SaeCtlCvMax{}. This is
Proposition~\ref{prop:leverage} in action rather than a coincidence: an i.i.d.\ code concentrates its
leverage tightly around $d/F$, and a code with equalised leverage sits at the floor by the identity.

The consequence for how such ratios are read is direct. A reported attainment ratio near one, under a
pseudoinverse readout, is consistent with a code that has learned nothing at all. It becomes evidence
about training only when compared against a shape-matched baseline, and that comparison is the first
thing we could not find in the work we build on.

\subsection{No released dictionary we measured is at the floor}

All \SaeDictsAboveControl{} of the \SaeDicts{} dictionaries sit above their own control, at $\Rgeom$
from \SaeDictRgeomMin{} to \SaeDictRgeomMax{}. There is no exception across two decoder sites
(residual stream and MLP), three training recipes (TopK, JumpReLU-style and a hybrid model's own
projection), five model families, and a twentyfold range of load.

The direction is worth pausing on, because it is the opposite of the natural guess. One might expect
a dictionary trained to represent a model's features to be at least as well arranged for linear
readout as an unstructured code. It is not: a random code of the same shape carries strictly less
cross-talk under the best linear readout each admits. Whatever these dictionaries optimise --- sparse
reconstruction of activations, under an $\ell_0$ or TopK constraint --- it is not the conditioning of
their own linear readout, and the two objectives are not aligned.

By the identity, the excess has a single source: leverage dispersion, \SaeDictCvMin{} to
\SaeDictCvMax{} against the controls' \SaeCtlCvMin{} to \SaeCtlCvMax{}. A released dictionary
allocates its directions unevenly with respect to the space it lives in, and a handful of features
with small $h_i$ dominate the harmonic sum that sets the ratio.

\paragraph{Depth, not width.} The excess is not uniform across a model. In every suite we measured,
mid-depth dictionaries sit furthest from the floor and the deepest layer sampled sits closest ---
$\Rgeom=\SaeDictRgeomMax{}$ at mid-depth against \SaeDictRgeomMin{} at the last layer of one family,
with the same ordering in the others. Whatever a practitioner is choosing when they pick a layer to
train a dictionary on, they are also choosing how well that dictionary can be read linearly, by a
margin far larger than any width effect below.

\subsection{A randomly initialised control shows where the excess comes from}

The obvious objection to the previous subsection is that training a TopK autoencoder might produce
uneven leverage on any input, in which case the excess would say nothing about the model. An i.i.d.\
control cannot answer this, because it does not share the training procedure.

EleutherAI released the pair that does. Two dictionaries trained by an identical recipe on identical
data, one on SmolLM2-135M and one on a \emph{randomly initialised} copy of it, at the same ten layers
and the same $64\times$ expansion. The randomly initialised arm returns to the floor:
$\Rgeom=\SaeRandInitRgeomMin$ to \SaeRandInitRgeomMax{} with dispersion \SaeRandInitCvMin{} to
\SaeRandInitCvMax{}, against the trained arm's excess in all \SaePairedLayers{} paired layers. Its
values are also nearly flat across depth, where the trained arm has clear structure.

Fitting a sparse autoencoder therefore does not, by itself, produce dispersed leverage. The excess
tracks what the model represents. This is the strongest causal statement the section supports, and it
rests on one model pair, which is the only such pair we are aware of --- a limitation we return to in
Section~\ref{sec:limitations}.

\subsection{Load and dispersion come apart}

Holding $d$ and $k$ fixed and varying only width, the ratio is flat while the dispersion grows:
$\Rgeom$ of \SaeLoadOneRgeom, \SaeLoadTwoRgeom{} and \SaeLoadThreeRgeom{} at loads of
\SaeLoadOneLoad$\times$, \SaeLoadTwoLoad$\times$ and \SaeLoadThreeLoad$\times$, while
$\mathrm{cv}(h)$ climbs from \SaeLoadOneCv{} through \SaeLoadTwoCv{} to \SaeLoadThreeCv{}.

Without the identity this reads as an anomaly. With it, it is a prediction: $\Rgeom$ is set by
$\sum_i h_i^{-1}$, a harmonic quantity dominated by the smallest leverage scores, whereas the
coefficient of variation is a second-moment summary sensitive to the upper tail. Extra spread at
higher load can therefore accumulate in a part of the distribution the ratio does not see. Across all
\SaeDicts{} dictionaries the excess is ranked almost perfectly by the harmonic-to-arithmetic gap of
the leverage and only well by the dispersion, which is the same statement measured rather than
derived.

Two things follow for practice. Widening a dictionary at fixed $d$ does not improve --- and does not
much worsen --- how well it can be read linearly. And the coefficient of variation, the obvious
summary to reach for, is the wrong one: it moves when the ratio does not.

\begin{table}[t]
\centering
\caption{The Interface Diagnostic on \SaeDicts{} released dictionaries, grouped by the axis each
block isolates, with every i.i.d.\ control matched to the shape and operating sparsity of the block
it judges. $s$ is the dictionary's own published operating sparsity. The
\emph{trained against randomly initialised} block is a matched intervention rather than an
observation: both arms come from the same recipe, data and layers, and differ only in whether the
underlying model was trained. Generated from \texttt{results/sae/derived/sae\_axes.json}; the SHA-256
of every checkpoint analysed is recorded there.}
\label{tab:sae}
\small
% Generated by scripts/make_tables.py from results/. Do not edit by hand.
\begin{tabular}{llrrrrrr}
\toprule
code & & $d$ & $F$ & $F/d$ & $s$ & $R_{\mathrm{geom}}$ & $\mathrm{cv}(h)$ \\
\midrule
\multicolumn{8}{l}{\itshape Scale: $d$ varies, load fixed at $16\times$} \\
Qwen3-1.7B layer1 &  & 2048 & 32768 & 16.0 & 50 & 1.1371 & 0.356 \\
Qwen3-1.7B layer18 &  & 2048 & 32768 & 16.0 & 50 & 1.1218 & 0.274 \\
Qwen3-1.7B layer27 &  & 2048 & 32768 & 16.0 & 50 & 1.0055 & 0.102 \\
Qwen3-1.7B layer9 &  & 2048 & 32768 & 16.0 & 50 & 1.1849 & 0.331 \\
Qwen3-8B L17 &  & 4096 & 65536 & 16.0 & 50 & 1.2018 & 0.352 \\
Qwen3-8B L30 &  & 4096 & 65536 & 16.0 & 50 & 1.0178 & 0.126 \\
Qwen3-8B \_L4 &  & 4096 & 65536 & 16.0 & 50 & 1.0812 & 0.262 \\
Qwen3.5-2B layer1 &  & 2048 & 32768 & 16.0 & 50 & 1.0336 & 0.180 \\
Qwen3.5-2B layer16 &  & 2048 & 32768 & 16.0 & 50 & 1.0798 & 0.259 \\
Qwen3.5-2B layer23 &  & 2048 & 32768 & 16.0 & 50 & 1.0306 & 0.187 \\
Qwen3.5-2B layer8 &  & 2048 & 32768 & 16.0 & 50 & 1.1704 & 0.356 \\
i.i.d. 2048x32768 & control & 2048 & 32768 & 16.0 & 50 & 1.0001 & 0.008 \\
i.i.d. 4096x65536 & control & 4096 & 65536 & 16.0 & 50 & 1.0000 & 0.005 \\
\addlinespace[3pt]
\multicolumn{8}{l}{\itshape Load: $F/d$ varies, $d$ and $k$ fixed} \\
Pythia-160m 4k &  & 768 & 4084 & 5.3 & 64 & 1.0674 & 0.276 \\
Pythia-160m 32k &  & 768 & 32768 & 42.7 & 64 & 1.0749 & 0.404 \\
Pythia-160m 65k &  & 768 & 65536 & 85.3 & 64 & 1.0734 & 0.469 \\
i.i.d. 768x4084 & control & 768 & 4084 & 5.3 & 64 & 1.0005 & 0.020 \\
i.i.d. 768x32768 & control & 768 & 32768 & 42.7 & 64 & 1.0001 & 0.008 \\
i.i.d. 768x65536 & control & 768 & 65536 & 85.3 & 64 & 1.0000 & 0.005 \\
\addlinespace[3pt]
\multicolumn{8}{l}{\itshape Trained model against a randomly initialised one} \\
SmolLM2 random L0 &  & 576 & 36864 & 64.0 & 64 & 1.0038 & 0.060 \\
SmolLM2 random L12 &  & 576 & 36864 & 64.0 & 64 & 1.0063 & 0.077 \\
SmolLM2 random L15 &  & 576 & 36864 & 64.0 & 64 & 1.0063 & 0.077 \\
SmolLM2 random L18 &  & 576 & 36864 & 64.0 & 64 & 1.0064 & 0.078 \\
SmolLM2 random L21 &  & 576 & 36864 & 64.0 & 64 & 1.0063 & 0.077 \\
SmolLM2 random L24 &  & 576 & 36864 & 64.0 & 64 & 1.0065 & 0.078 \\
SmolLM2 random L27 &  & 576 & 36864 & 64.0 & 64 & 1.0065 & 0.078 \\
SmolLM2 random L3 &  & 576 & 36864 & 64.0 & 64 & 1.0045 & 0.066 \\
SmolLM2 random L6 &  & 576 & 36864 & 64.0 & 64 & 1.0057 & 0.074 \\
SmolLM2 random L9 &  & 576 & 36864 & 64.0 & 64 & 1.0061 & 0.076 \\
SmolLM2 trained L0 &  & 576 & 36864 & 64.0 & 64 & 1.0512 & 0.307 \\
SmolLM2 trained L12 &  & 576 & 36864 & 64.0 & 64 & 1.0302 & 0.248 \\
SmolLM2 trained L15 &  & 576 & 36864 & 64.0 & 64 & 1.1160 & 0.351 \\
SmolLM2 trained L18 &  & 576 & 36864 & 64.0 & 64 & 1.1351 & 0.397 \\
SmolLM2 trained L21 &  & 576 & 36864 & 64.0 & 64 & 1.0253 & 0.198 \\
SmolLM2 trained L24 &  & 576 & 36864 & 64.0 & 64 & 1.0355 & 0.279 \\
SmolLM2 trained L27 &  & 576 & 36864 & 64.0 & 64 & 1.0319 & 0.211 \\
SmolLM2 trained L3 &  & 576 & 36864 & 64.0 & 64 & 1.0105 & 0.177 \\
SmolLM2 trained L6 &  & 576 & 36864 & 64.0 & 64 & 1.0216 & 0.190 \\
SmolLM2 trained L9 &  & 576 & 36864 & 64.0 & 64 & 1.0207 & 0.198 \\
i.i.d. 576x36864 & control & 576 & 36864 & 64.0 & 64 & 1.0001 & 0.007 \\
\addlinespace[3pt]
\multicolumn{8}{l}{\itshape Hybrid convolution/attention MLP code} \\
LFM2.5-350M L0 (conv) &  & 1024 & 4608 & 4.5 & 64 & 1.0110 & 0.101 \\
LFM2.5-350M L1 (conv) &  & 1024 & 4608 & 4.5 & 64 & 1.0453 & 0.247 \\
LFM2.5-350M L10 (attn) &  & 1024 & 4608 & 4.5 & 64 & 1.0510 & 0.239 \\
LFM2.5-350M L11 (conv) &  & 1024 & 4608 & 4.5 & 64 & 1.0330 & 0.204 \\
LFM2.5-350M L12 (attn) &  & 1024 & 4608 & 4.5 & 64 & 1.0119 & 0.122 \\
LFM2.5-350M L13 (conv) &  & 1024 & 4608 & 4.5 & 64 & 1.0247 & 0.196 \\
LFM2.5-350M L14 (attn) &  & 1024 & 4608 & 4.5 & 64 & 1.0083 & 0.090 \\
LFM2.5-350M L15 (conv) &  & 1024 & 4608 & 4.5 & 64 & 1.0240 & 0.167 \\
LFM2.5-350M L2 (attn) &  & 1024 & 4608 & 4.5 & 64 & 1.0179 & 0.133 \\
LFM2.5-350M L3 (conv) &  & 1024 & 4608 & 4.5 & 64 & 1.0183 & 0.158 \\
LFM2.5-350M L4 (conv) &  & 1024 & 4608 & 4.5 & 64 & 1.0605 & 0.262 \\
LFM2.5-350M L5 (attn) &  & 1024 & 4608 & 4.5 & 64 & 1.0267 & 0.154 \\
LFM2.5-350M L6 (conv) &  & 1024 & 4608 & 4.5 & 64 & 1.0197 & 0.149 \\
LFM2.5-350M L7 (conv) &  & 1024 & 4608 & 4.5 & 64 & 1.0528 & 0.242 \\
LFM2.5-350M L8 (attn) &  & 1024 & 4608 & 4.5 & 64 & 1.0330 & 0.168 \\
LFM2.5-350M L9 (conv) &  & 1024 & 4608 & 4.5 & 64 & 1.0626 & 0.297 \\
i.i.d. 1024x4608 & control & 1024 & 4608 & 4.5 & 64 & 1.0004 & 0.018 \\
\addlinespace[3pt]
\bottomrule
\end{tabular}

\end{table}
